\section{Preliminaries and Motivation}
\subsection{Problem Setting}
\noindent\textbf{Task Setup.}
We consider a long-horizon conversational setting in which an agent processes a stream of dialogue chunks
$\mathcal{C}=\{c_1,\ldots,c_T\}$ over time. 
The stream is further organized into sessions $\mathcal{S}=\{s_1,\ldots,s_N\}$,where each session $s_\tau$ consists of one or more consecutive chunks corresponding to a coherent interaction episode.
At each step $t$, the agent maintains an external memory bank $M_t$ composed of structured entries
(e.g., text, timestamp, source, and speaker metadata), which is updated as new conversational information arrives.
After processing the full stream $\mathcal{C}$, the agent is evaluated on a set of questions $Q$.
For each query $q \in Q$, it retrieves evidence $E(q)$ from $M_T$ and outputs an answer $\hat{y}(q)$.
Our goal is to design an agent $\pi$ that maximizes answer accuracy by jointly improving memory construction and retrieval.

% After observing the full chunks $\mathcal{C}$, the agent is evaluated on a set of queries ${Q}$: for each $q\in{Q}$, it retrieves evidence $E(q)$ from $M_T$ and outputs an answer ${y}(q)$. Our objective is to learn $\pi$ that maximizes answer accuracy by jointly optimizing memory construction and retrieval. 

\noindent\textbf{Challenges.}
This setting is challenging because success depends on both memory construction and memory retrieval.
During \emph{construction}, the agent must decide what to write, update, merge, or discard when a new chunk arrives.
During \emph{retrieval and answering}, it must identify the right evidence from memory under ambiguity, temporal dependencies, and incomplete or underspecified initial queries.
The challenge is therefore not merely to improve answer generation, but to maintain a useful memory bank and retrieve the right evidence under bounded memory and retrieval budgets.


\subsection{Memory Cycle Effect as a Design Lens}
\label{ssec:cycle_effect}
The above challenges suggest that long-term memory should not be viewed as a linear pipeline of isolated modules.
Instead, we adopt the \textbf{memory cycle effect}~\cite{zhang2025learn} as a design lens for analyzing long-term memory systems.
Under this view, memory forms a closed loop with three tightly coupled phases: \emph{construction}, \emph{retrieval}, and \emph{utilization}.
Construction determines what information enters the memory bank and how it is organized; retrieval determines what stored information is surfaced as evidence; and utilization reveals whether the retrieved evidence is sufficient for downstream answering.

This perspective highlights two dependencies. First, there is a \emph{forward dependency}: construction constrains retrieval, and retrieval in turn constrains utilization.
A poorly constructed memory bank may omit important details, retain redundant entries, or leave conflicts unresolved, all of which degrade downstream retrieval quality. 
Second, there is a \emph{backward dependency}: utilization outcomes expose deficiencies in upstream memory operations, since answering failures may stem from earlier storage omissions, unresolved contradictions, or poorly targeted retrieval. 
As a result, the utility of memory operations is often sparse and delayed, making isolated optimization of memory modules fundamentally suboptimal.

Together, these dependencies suggest that long-term memory should be studied as a coupled cycle rather than independent storage and retrieval components.
This motivates the need for mechanisms that explicitly coordinate forward memory execution and propagate utilization feedback backward to improve future memory decisions.




\begin{table}[t]
\centering
% \scriptsize
% \small

% \vspace{-1em}
\caption{Preliminary analysis results (\%) in LoCoMo dataset, GPT-4o-mini is the backbone LLM. 
% \zhiwei{since the performance gap for F1 is not that significant, could we report variance}
} 
% \vspace{-1em}
\label{tab:prior_study}
\scalebox{0.95}{\begin{tabular}{lccc}
\toprule
\textbf{Method} & \textbf{F1} & \textbf{B1} & \textbf{ACC} \\%& 
     \midrule
Static Baseline & 22.64 & 17.24 & 52.60 \\
Unguided Active & 23.49 & \textbf{18.36} & 54.60\\
Strategic Active & \textbf{24.78} & 17.73 & \textbf{59.21}\\
\bottomrule
\end{tabular}}
% \vspace{-1.5em}
\end{table}


\subsection{Motivating Analysis: Strategic Blindness}
\label{ssec:motivate_study}
% \suhang{the logic here is quite strange. We were talking about ``the need for xxx'' in the above paragraph. Then we suddenly talk about ``recent agents have moved beyond fully passive to xxx''. I do not get the connection or the reason why we want to discuss this. Are we trying to do experiments to further demonstrate ``the need for mechanisms that xxx'', or are we trying to show another issue? Could we add some transition sentences here?} 
The analysis above motivates coordination across the memory cycle, 
but do existing active memory agents achieve this in practice?
Recent agents~\cite{fang2025lightmem,xu2025mem} have moved beyond 
fully passive memory by introducing active updates or iterative 
retrieval. However, most still operate in a largely reactive 
manner: they trigger operations based on local context or immediate similarity signals rather than an explicit global strategy. 
% Recent memory-augmented agents~\cite{fang2025lightmem,xu2025mem} have moved beyond fully passive memory by introducing active updates or iterative retrieval.
% However, most of them still operate in a largely reactive manner: they trigger operations based on local context or immediate similarity signals rather than an explicit global strategy. 
We characterize this limitation as \emph{strategic blindness}: the agent has the \emph{hands} to edit memory and issue retrieval queries, but lacks the \emph{brain} to coordinate these actions across the full memory cycle.
% We identify two concrete manifestations:
This manifests as:
(i)~\emph{Myopic Construction}: construction decisions are driven 
by local context rather than downstream utility. The agent 
indiscriminately appends, overwrites, or ignores information, 
leaving redundancy and conflicts unresolved.
(ii)~\emph{Aimless Retrieval}: when the initial query is incomplete 
or semantically mismatched with stored memory, one-shot retrieval 
or shallow rewrites fail to surface the required evidence. Without 
strategic guidance, successive queries do not narrow the information gap.

\noindent\textbf{Setup.} 
To empirically validate this diagnosis, we conduct a preliminary study on a subset of LoCoMo~\cite{maharana2024evaluating}, focusing on reasoning-intensive queries by excluding adversarial samples.
We compare three progressively stronger baselines using GPT-4o-mini~\cite{gpt4ocard} as the backbone:
(i) \textit{Static}, which performs memory construction followed by one-shot top-$30$ retrieval;
(ii) \textit{Unguided Active}, which adds iterative query rewriting without strategic guidance; and
(iii) \textit{Strategic Active}, which introduces a planner to guide both construction and retrieval.
We report token-level F1, BLEU-1 (B1), and LLM-as-a-Judge accuracy (ACC).
More evaluation details are provided in Appendix~\ref{appendix:preliminary_evaluation_details}.


% \noindent\textbf{Empirical analysis.} 
% Table~\ref{tab:prior_study} reveals two findings:
% (i)~\emph{Refinement provides capability:} the \textit{Unguided Active} variant (54.6\% ACC) outperforms the \textit{Static} variant (52.6\%), confirming that iterative refinement helps bridge the semantic gap missed by one-shot retrieval;
% (ii)~\emph{Reasoning provides control:} the \textit{Strategic Active} variant achieves a larger leap to 59.2\% ACC, substantially outperforming both weaker variants.
% Case studies in Appendix~\ref{appendix:preliminary_case_studies} further show that unguided memory storage and retrieval often induce reasoning drift and degrade utility.
% These findings suggest that active memory operations alone are insufficient: explicit strategic reasoning is needed to guide both construction and retrieval. \suhang{how does this experiment verify Myopic
% Construction and Aimless Retrieval? Could you explicitly connect the experiment design with these two points?}

\noindent\textbf{Empirical analysis.} Table~\ref{tab:prior_study} reveals two findings: (i)~\emph{Refinement provides capability:} {Unguided Active} (54.6\% Acc) outperforms {Static} (52.6\%), confirming that one-shot retrieval often fails to surface the required evidence when the initial query is incomplete or mismatched with memory, which directly reflects \emph{Aimless Retrieval}. (ii)~\emph{Reasoning provides control:} {Strategic Active} achieves a larger leap to 59.2\% Acc. Since it shares the same active operators as {Unguided Active}, this gap reflects the value of explicit strategic guidance in addressing both \emph{Aimless Retrieval} and \emph{Myopic Construction}. Case studies in Appendix~\ref{appendix:preliminary_case_studies} further illustrate both pathologies with concrete examples of redundant entries and retrieval drift. 
These findings suggest that active memory operations alone are insufficient: explicit strategic reasoning is needed to guide both construction and retrieval. 
% These findings motivate the design of \method, where a Meta-Thinker provides explicit coordination for both construction and retrieval.
% These findings motivate the forward-path design of \method, in which a Meta-Thinker provides explicit strategic reasoning for both construction and retrieval.



% \noindent\textbf{Empirical analysis.} 
% Table~\ref{tab:prior_study} reveals two findings:
% (i)~\emph{Refinement provides capability:} the \textit{Unguided Active} variant (54.6\% J) outperforms the \textit{Static} variant (52.6\%), confirming that iterative refinement helps bridge the semantic gap that one-shot retrieval misses;
% (ii)~\emph{Reasoning provides control:} the \textit{Strategic Active} variant achieves a larger leap to 59.2\% J, significantly outperforming both other variants.
% Case studies in Appendix~\ref{appendix:preliminary_case_studies} further show that unguided memory storage and retrieval often induces reasoning drift and degrades utility.
% This demonstrates that explicit reasoning provides the \emph{control} needed to filter noise during construction and target evidence during retrieval.

% Table~\ref{tab:prior_study} reveals two clear findings that support our hypothesis:
% (i) \emph{Refinement provides capability:} the \textit{unguided active} variant ($54.6\%$) outperforms the \textit{static} variant ($52.6\%$), confirming that iterative active refinement helps bridge the semantic gap that static retrieval misses;
% (ii) \emph{Reasoning provides control:} The \textit{strategic baseline} creates a larger leap to ${59.2\%}$, significantly outperforming the other variants. In addition, our case studies in Appendix~\ref{appendix:preliminary_case_studies} further reveal that unguided memory storage and retrieval often induces reasoning drift and degrades utility. This demonstrates that explicit reasoning provides the \textit{control}, filtering noise during construction and targeting evidence during retrieval.

% \zhiwei{since the experiment results would motivate the core design of our framework, could we have more results using different benchmarks, base models, or even planning agent}
